\documentclass{beamer}

% The legendary old-school blue Beamer theme
\usetheme{Madrid}
\usecolortheme{default}

\title{Debt Simplification Algorithm}
\subtitle{Graph Theory Application}
\author{Algorithm Overview}
\date{\today}

\begin{document}

% --- Frame 1: Title Page ---
\begin{frame}
    \titlepage
\end{frame}

% --- Frame 2: The Task & Model ---
\begin{frame}{Task \& Model Definition}
    \textbf{Graph Representation:}
    \begin{itemize}
        \item People are represented as \textbf{vertices}.
        \item Debts are represented as \textbf{directed, weighted edges}.
    \end{itemize}

    \vspace{0.5cm}
    \textbf{The Objective:}
    \begin{itemize}
        \item Given an initial, complex directed graph $X$.
        \item Create a new, simplified directed graph $X'$.
    \end{itemize}
\end{frame}

% --- Frame 3: The Core Invariant ---
\begin{frame}{Necessary Properties of $X'$}
    \textbf{The Core Rule:}
    For every person (``guy") in the graph, their net balance must remain exactly the same.

    \vspace{0.5cm}
    \begin{block}{Net Balance Equation}
        Let $Net(v) = IN(v) - OUT(v)$.
    \end{block}

    \vspace{0.3cm}
    Therefore, for all vertices $v$:
    \begin{itemize}
        \item $Net(v)$ in graph $X$
        \item \textbf{must equal} $Net(v)$ in graph $X'$
    \end{itemize}
\end{frame}

% --- Frame 4: Initialization ---
\begin{frame}{Constructing $X'$: Initialization}
    \begin{enumerate}
        \item Initially, the new graph $X'$ has \textbf{no edges}.
        \vspace{0.2cm}
        \item Calculate the remaining balance required for everyone:
        \[ \Delta Net(v) = Net_{X}(v) - Net_{X'}(v) \]
        \item Divide people into two distinct groups:
        \begin{itemize}
            \item \textbf{Creditors:} $\Delta Net(v) > 0$
            \item \textbf{Debtors:} $\Delta Net(v) < 0$
        \end{itemize}
        \vspace{0.2cm}
        \item[] \textit{Note:} We completely ignore people with $\Delta Net(v) = 0$, because their required property is already successfully satisfied.
    \end{enumerate}
\end{frame}

% --- Frame 5: Edge Creation ---
\begin{frame}{Constructing $X'$: Edge Creation}
    \textbf{Iterative Simplification Step:}

    \vspace{0.3cm}
    \begin{itemize}
        \item Choose two people, $A$ and $B$, from \textbf{different groups}.
        \vspace{0.2cm}
        \item Create a directed edge between them in $X'$ with a weight of:
        \[ \text{Weight} = \min(|\Delta Net(A)|, |\Delta Net(B)|) \]
        \vspace{0.2cm}
        \item Subtract this weight from their respective balances.
        \vspace{0.2cm}
        \item \textbf{Result:} Because we subtract the minimum of the two absolute values, at least one person's $\Delta Net$ becomes exactly zero, settling them for good!
    \end{itemize}
\end{frame}

% --- Frame 6: Complexity & Edge Count ---
\begin{frame}{Constructing $X'$: Number of Edges}
    \textbf{Upper Bound on Edges in $X'$:}

    \vspace{0.2cm}
    The maximum number of edges in the simplified graph is \textbf{$N - 1$}, where $N$ is the total number of people with a non-zero initial net balance.

    \vspace{0.5cm}
    \textbf{Why?}
    \begin{itemize}
        \item Every new edge added is weighted at $\min(|\Delta Net(A)|, |\Delta Net(B)|)$.
        \item This guarantees at least one person's $\Delta Net$ reaches $0$ per edge.
        \item Therefore, each edge permanently resolves and removes at least one person from the remaining pool.
        \item The final edge added will balance out the final creditor and final debtor simultaneously.
    \end{itemize}
\end{frame}

\end{document}